The proliferation of Large Language Models (LLMs) has sparked an arms race between AI-generated text detectors and evasion techniques. While existing research focus on complex algorithmic optimizations or direct prompting instructions, we investigate the "Bullying Paradox": the phenomenon where iterative negative feedback (e.g., rejecting output as "too AI-sounding") significantly outperforms a single positive instruction to "sound like a human." We systematically evaluate this effect using a conversational feedback loop across five rounds of rejection. Our results show that iterative "bullying" achieves a \roundfiveacc average "Real" probability score, nearly doubling the \singleshotacc score of single-shot instructions. Through linguistic analysis, we identify that this performance gap is driven by a qualitative regime shift in the model's generation, characterized by a 20.5\% increase in perplexity and a substantial rebound in sentence length variation (burstiness). Our findings suggest that current detectors are highly vulnerable to simple conversational dynamics that steer models into previously inaccessible, human-like latent regions.
